\section{From finite tests to an actual null map}\label{sec:upper-nullity}

We now bound the exponent of the upper local unit by proving nullity
of a specified tensor-power map.

\begin{proposition}[The upper null map and its finite layers]
\label{wup:upper-null}
Let \(p\geq5\) be prime and put \(h=p\). In \(\mathcal C_h\), let
\(I_h=\fib(B_h\to E_h)\). The actual map
\[
 e_m(B_h):I_h^{\otimes_h m}\longrightarrow B_h,
 \qquad m=h^2+2=p^2+2,
\]
is null. Thus \(\expn_{E_h}(B_h)\leq m\), and \(B_h\) is a retract
of an object with \(m\) successive layers
\[
 E_h\otimes_h I_h^{\otimes_h i},\qquad 0\leq i<m.
\]
\end{proposition}

The coefficient line will make this map zero after every finite ordinary
precomposition, once some finite exponent is known. The following
property of its target then turns that tested vanishing into nullity.
Its relation to the Brown--Comenetz phantom criterion is
\cite[Proposition~4.11]{christensen-strickland1998}.

\begin{lemma}[Ordinary phantoms into a local sphere]\label{wup:phantom}
Let \(a\geq1\). An ordinary map \(f:U\to B_a\) of \(p\)-local
spectra is null if its precomposition with every map from a finite
ordinary \(p\)-local spectrum to \(U\) is null.
\end{lemma}

\begin{proof}
Let \(M_aS=\fib(L_aS\to L_{a-1}S)\), and let \(I_{\mathrm{BC}}\)
be the ordinary Brown--Comenetz dualizing spectrum for the injective
group \(\Q/\Z_{(p)}\). Gross--Hopkins duality, announced in
\cite[Theorem~6]{gross-hopkins1994} and proved in the form used here
in \cite[Theorem~10.2(b),(e)]{hs1999}, says that
\[
 D_a=F(M_aS,I_{\mathrm{BC}})
\]
is \(K(a)\)-local and invertible for \(\otimes_a\). The
identification with this monochromatic object uses
\cite[Theorem~6.19]{hs1999}.

For a local target \(Z\), the ordinary spectrum \(F(R,Z)\) is
local. A map into it from a \(K(a)\)-acyclic spectrum \(V\)
corresponds to a map \(V\wedge R\to Z\), and \(V\wedge R\) is
still acyclic. Hence ordinary function spectra compute the local
internal Hom, as in \cite[Theorem~7.1]{hs1999}. Invertibility of
\(D_a\), followed by ordinary closed adjunction, gives equivalences
of ordinary spectra
\[
 B_a\simeq F(D_a,D_a)
       \simeq F(D_a\wedge M_aS,I_{\mathrm{BC}}).
\]
The displayed smash is ordinary. Put \(C=D_a\wedge M_aS\).
Brown--Comenetz representability now gives, naturally in \(U\),
\[
 [U,B_a]\cong
 \operatorname{Hom}_{\Z}\bigl(\pi_0(U\wedge C),\Q/\Z_{(p)}\bigr).
\]

Write \(U\) as a filtered homotopy colimit of finite ordinary
\(p\)-local spectra \(U_i\). Ordinary smash with \(C\) and ordinary
homotopy groups preserve this colimit, so
\[
 \pi_0(U\wedge C)=\varinjlim_i\pi_0(U_i\wedge C).
\]
The functional corresponding to \(f\) vanishes on every term by
the assumed finite precompositions. It therefore vanishes on the
colimit, and \(f=0\). This argument tests by ordinary finite
spectra and uses no interchange of \(K(a)\)-localization with
this filtered colimit.
\end{proof}

\begin{proof}[Proof of Proposition~\ref{wup:upper-null}]
First obtain some finite exponent for the actual local unit,
independently of any finite test. Devinatz--Hopkins prove that
every \(K(h)\)-local spectrum belongs to the thick tensor ideal
generated by \(E_h\) in \(\mathcal C_h\)
\cite[Appendix~A, Proposition~A.3]{dh2004}. This is the local
nilpotence consequence of Hopkins--Ravenel; its descendability
form is \cite[Proposition~10.10]{mathew2016}. The objects of finite
\(E_h\)-exponent form a thick tensor ideal containing \(E_h\):
modules have exponent at most one, tensoring preserves a bound,
and cofiber sequences and retracts preserve finite bounds by
Lemma~\ref{lem:exponent-calculus}. It follows that \(B_h\) has
some finite exponent \(c\). This is one bound on an actual map
before any test is chosen; the nilpotence theorem supplies no
numerical value needed here.

Apply the full local-unit Amitsur construction in \(\mathcal C_h\)
to \(B_h\), and then the exact ordinary functor \(F(V_0,-)\) for a
finite ordinary \(V_0\). Proposition~\ref{wup:cochains} identifies
its second page with
\[
 H^s_{\cts}\bigl(\G_h,\pi_q(E_h\wedge D V_0)\bigr).
\]
Proposition~\ref{prop:cohomology} makes this page zero for
\(s\geq h^2+2=m\). The tested-image assertion of
Proposition~\ref{prop:pages}(ii) applies using the already known
finite exponent \(c\). It concerns the full finite partial
totalizations and their finite fiber cubes, and gives
\[
 [V_0,I_h^{\otimes_h m}]\longrightarrow[V_0,B_h]
       \quad\text{zero for every finite ordinary }V_0.
\]
Thus the actual arrow \(e_m(B_h)\), regarded in ordinary spectra,
vanishes after every finite ordinary precomposition.
Lemma~\ref{wup:phantom} makes it null as an ordinary map and
hence null in the full subcategory \(\mathcal C_h\).

For clarity, all layers needed next occur at finite stages. Set
\(I_h^{\otimes_h0}=B_h\), and for \(0\leq j\leq m\) put
\[
 F_j=\cofib\bigl(
 I_h^{\otimes_h m}\longrightarrow I_h^{\otimes_h(m-j)}
 \bigr).
\]
The consecutive maps form a finite filtration
\(0=F_0\to F_1\to\cdots\to F_m\).
The cofiber of \(I_h^{\otimes_h(i+1)}\to I_h^{\otimes_h i}\)
is \(E_h\otimes_h I_h^{\otimes_h i}\). The cofiber comparison
for consecutive composites therefore identifies
\[
 \cofib(F_j\to F_{j+1})
 \simeq E_h\otimes_h I_h^{\otimes_h(m-j-1)}
       \quad(0\leq j<m).
\]
Finally \(F_m=\cofib(e_m(B_h))\). A nullhomotopy of \(e_m(B_h)\)
makes \(B_h\) a retract of this cofiber. This proves both assertions.
\end{proof}
